% Lösungen Einheit 1 von 5 – Gleichungen mit zwei Variablen und grafisches Lösen (zusammenbau v0.1)
\einheitenkopf{Einheit 1 von 5 · Gleichungen mit zwei Variablen und grafisches Lösen}

\erg{13}{a) von beiden \quad b) $y$: $5$, $4$, $3$, $2$ \quad c) $y = -2x + 6$ \quad d) $y = x + 1$: Gerade durch $(0|1)$ und $(2|3)$ \quad e) Schnittpunkt $(2|3)$ \quad f) I: $3 + 5 = 8$ (wA); II: $1 - 5 = -4$ (wA); ja, $(1|5)$ ist die Lösung \quad g) $(2|1)$ \quad h) keine \quad i) $x = 3$, $y = 5$ \quad j) $x = 6$ Zweierzelte, $y = 8$ Viererzelte; Probe: $6 + 8 = 14$, $12 + 32 = 44$. Keine andere Möglichkeit: Tauscht man ein Zweierzelt gegen ein Viererzelt, kommen $2$ Plätze dazu – die Platzzahl ändert sich bei jedem Tausch, nur eine Aufteilung ergibt $44$.}
\erg{14}{a) ein Schnittpunkt \quad b) Gerade $y = x - 2$; Lösung mit $y = 3$: $x = 5$ \quad c) I: $2 \cdot 1 + 3 = 5$ (wA); II: $1 - 3 \cdot 3 = -8$ (wA). II ist kein Vielfaches von I, die Gleichungen sind nicht äquivalent: Ihre Geraden sind verschieden und nicht parallel, sie schneiden sich in genau einem Punkt – es gibt keine weitere Lösung.}
\erg{15}{a) Probe in II: $-1 + 3 = 2$, nicht $1$ (fA). Der Schnittpunkt liegt zwischen Gitterpunkten: $2x - 1 = -x + 3$ ergibt $x = \frac{4}{3}$, $y = \frac{5}{3}$. \quad b) $3x$ wird auf beiden Seiten abgezogen: $y = -3x + 5$. \quad c) Aus I folgt $y = 3x + 2$: Die Steigung ist $3$, nicht $-3$. \quad d) Jede Gerade besteht aus allen Lösungen ihrer Gleichung. Der Schnittpunkt liegt auf beiden Geraden, also löst er beide Gleichungen. \quad e) Beide Geraden haben die Steigung $3$, aber verschiedene Achsenabschnitte: Sie sind parallel und haben keinen gemeinsamen Punkt. \quad f) Jede Gleichung beschreibt eine Gerade. Sind die Gleichungen nicht äquivalent, sind die Geraden verschieden; zwei verschiedene Geraden schneiden sich in höchstens einem Punkt oder sind parallel.}
